\chapter{Überblick}

Der League Simulator schätzt für zehn deutsche Fußballligen, mit welcher
Wahrscheinlichkeit jede Mannschaft am Saisonende auf welchem Tabellenplatz
steht. Er rechnet nach jedem beendeten Spiel neu und veröffentlicht das
Ergebnis als statische Website unter
\href{https://fussball.csdatascience.de}{fussball.csdatascience.de}, die unter
dem Namen des Blogs \enquote{30~Punkte} läuft.

Abgedeckt sind die Bundesliga, die 2.~Bundesliga, die 3.~Liga, die fünf
Regionalligen (Nord, Nordost, West, SüdWest, Bayern) sowie die
Frauen-Bundesliga und die 2.~Frauen-Bundesliga.

\section*{Das Modell in drei Sätzen}

\begin{enumerate}
  \item Jede Mannschaft hat eine \emph{Stärkeschätzung}, den ELO-Wert. Er steigt,
        wenn sie besser abschneidet, als das Modell erwartet hat, und fällt,
        wenn sie schlechter abschneidet (Kapitel~\ref{kap:elo}).
  \item Aus dem ELO-Unterschied zweier Mannschaften folgt für jede eine
        Torerwartung. Die Tore eines Spiels werden als unabhängige
        Poisson-Zufallszahlen gezogen (Kapitel~\ref{kap:tormodell}).
  \item Der Rest der Saison wird zehntausendmal durchgespielt. Der Anteil der
        Durchläufe, in denen eine Mannschaft auf einem Platz landet, ist die
        angezeigte Wahrscheinlichkeit (Kapitel~\ref{kap:simulation}).
\end{enumerate}

Alles Weitere ist Feinarbeit: Regeln für Auf- und Abstieg zwischen den Ligen
(Kapitel~\ref{kap:regionalliga}), die Herkunft der Startwerte
(Kapitel~\ref{kap:startwerte}) und eine ehrliche Liste dessen, was das Modell
nicht kann (Kapitel~\ref{kap:grenzen}).

\section*{Aufbau dieses Handbuchs}

\begin{description}[font=\sffamily\bfseries\small]
  \item[Teil I, Methodik] Alle Formeln, so wie der Code sie rechnet,
        mit Konstanten, Rechenbeispielen und Quellverweisen.
  \item[Teil II, Website] Die 22 Seiten und jedes Element darauf, von der
        Heatmap bis zum ELO-Verlaufsdiagramm, dazu das Gestaltungssystem.
  \item[Teil III, Technik] Datenfluss, Scheduler, Rust-Simulationsserver,
        Rendering, Container, CI und Betrieb.
\end{description}

\begin{merke}[Zur Genauigkeit]
Die Formeln in diesem Handbuch sind aus dem Quelltext abgelesen und nicht
aus Lehrbüchern übernommen. Wo der Code von der üblichen Darstellung
abweicht, etwa beim Tordifferenz-Multiplikator oder bei der Begrenzung der
ELO-Differenz, gilt der Code. Verweise wie
\datei{league-simulator-rust/src/elo/mod.rs:5} nennen Datei und Zeile auf dem
Stand \texttt{a41890a}. Alle Rechenbeispiele sind nachgerechnet.
\end{merke}

\section*{Die Modellkonstanten}

Zwei Parametergruppen bestimmen jede Prognose. Sie stehen
\emph{ausschließlich} im Rust-Server
(\datei{league-simulator-rust/src/models/mod.rs:124--130}). R sendet einen Wert
nur, wenn eine Liga bewusst abweicht, und das betrifft nur das Tormodell der
Frauen-Ligen (ADR~0002, ADR~0004). Ein Wächtertest im R-Testsatz sorgt dafür,
dass keine dieser Zahlen im R- oder JavaScript-Code oder im Seitentext
auftaucht.

\begin{center}
\begin{tabular}{@{}llll@{}}
\toprule
Symbol & Name im Code & Wert & Gilt für \\
\midrule
$k$   & \code{mod\_factor}      & 20 & alle Ligen \\
$H$   & \code{home\_advantage}  & 40 & alle Ligen \\
$N$   & \code{iterations}       & \num{10000} & alle Ligen \\
$s$   & \code{tore\_slope}      & 0{,}0017854953143549 & Herren-Ligen \\
$c$   & \code{tore\_intercept}  & 1{,}3218390804597700 & Herren-Ligen \\
$s_F$ & \code{tore\_slope}      & 0{,}0024058833 & Frauen-Ligen (82, 1034) \\
$c_F$ & \code{tore\_intercept}  & 1{,}6527603153 & Frauen-Ligen (82, 1034) \\
\bottomrule
\end{tabular}
\end{center}
